% ============================================================
%  Machine Learning Project Report -- LaTeX edition, worked example
%  A FICTIONAL, generic sample used to exercise every template
%  feature. It is not a real submission and describes no real
%  person or organisation -- swap in your own content. Mirrors
%  example.typ so the two templates can be compared side by side.
%  Compile with: xelatex example.tex (run twice).
%
%  Every helper used below is documented -- LaTeX and Typst side by
%  side -- in docs/GUIDE.md under "Component reference".
% ============================================================
\documentclass[11pt,a4paper]{report}
\usepackage{utad}

\renewcommand{\utadTitle}{Image Classification with CNNs}
\renewcommand{\utadSubtitle}{Machine Learning Project}
\renewcommand{\utadDegree}{B.S. in Software Engineering, Minor in Data Science \& AI}
\renewcommand{\utadSubject}{Machine Learning}
\renewcommand{\utadYear}{4}
\renewcommand{\utadTeacher}{Teacher Name}
\renewcommand{\utadDate}{May 2026}
\setutadrunningtitle{ML Project Report}
\setutadauthor{Your Name}

\begin{document}
\utadtitlepage
\utadcontents

\chapter{Abstract}
\begin{utadnote}[Summary]
A convolutional neural network is trained to classify the Fashion-MNIST
dataset (28$\times$28 grayscale images, 10 clothing categories). Starting
from a plain baseline, the model is refined across eight tracked revisions;
the final network reaches competitive test accuracy, evaluated once on a
held-out set.
\end{utadnote}

\chapter{Introduction}

\section{Context and Motivation}
Image classification is a canonical supervised-learning task and a natural
vehicle for practicing the full machine-learning workflow: preparing data,
choosing and iterating on a model, tuning it against a validation set, and
reporting an honest estimate of generalization. This project takes
Fashion-MNIST --- a drop-in, harder replacement for the classic
handwritten-digit dataset --- and builds a convolutional classifier for it
from a simple baseline upward.

The task is a good fit for the degree's focus in that it combines software
engineering (a clean, reproducible training pipeline in Python) with the
data-science reasoning behind each modeling decision, so that every
accuracy gain can be attributed to a specific, justified change rather than
to undocumented trial and error.

\section{Objectives}
\begin{enumerate}
\item \textbf{General objective:} To apply and consolidate the
  machine-learning competencies acquired during the degree by building,
  training, and evaluating an image classifier end to end.
\item Iteratively refine a convolutional model, replacing a naive baseline
  with a tuned network whose every change is measured against a fixed
  validation split.
\item Apply sound evaluation practice --- a single, final test-set
  measurement --- so the reported accuracy is an unbiased estimate of
  generalization.
\item Document the modeling decisions and their measured effect clearly
  enough that the experiment is reproducible.
\end{enumerate}

\section{Report Structure}
This report is organized as follows: Chapter 3 describes the dataset, the
pipeline architecture, and the tools used. Chapter 4 details the modeling
work, organized by workstream. Appendix A provides a week-by-week
development log.

\chapter{Dataset and Approach}

\section{Dataset}
Fashion-MNIST consists of 70{,}000 grayscale images (60{,}000 train,
10{,}000 test) at 28$\times$28 resolution, in ten balanced classes
(t-shirt, trouser, pullover, dress, coat, sandal, shirt, sneaker, bag,
ankle boot). A fixed 10\% slice of the training set is held out as a
validation split; the official test set is reserved for a single final
measurement.

\section{Pipeline Architecture}
The classifier is organized as a small pipeline: preprocessing, a
convolutional feature extractor, and a dense classifier head. The
convolutional feature extractor (highlighted) was the focus of the
iteration work described in Chapter 4.

\begin{figure}[h]
\centering
\begin{utadorgchart}
\node[orgroot]{Classification Pipeline}
  child { node[orgnode]{Preprocessing}
    child { node[orgnode]{Normalize / augment} } }
  child { node[orgnode]{CNN Model}
    child { node[orghi]{Conv feature extractor} }
    child { node[orgnode]{Classifier head} } }
  child { node[orgnode]{Evaluation} };
\end{utadorgchart}
\caption{Architecture of the classification pipeline.}
\end{figure}

\section{Tools and Libraries}
{\utadtable
\begin{tabularx}{\textwidth}{@{}p{4cm}X@{}}
\tblheadrow \tblhead{Category} & \tblhead{Tools / Libraries} \\
Language & Python 3.11 \\
Modeling & PyTorch, torchvision \\
Numerics & NumPy \\
Plotting & Matplotlib \\
Experiment Tracking & Weights \& Biases (logging only) \\
Environment & Jupyter, Conda \\
Version Control & Git, GitLab (feature-branch workflow) \\
\end{tabularx}}

\chapter{Experiments}

The workstreams below are presented topically rather than strictly in
calendar order; data preparation and early modeling genuinely overlapped in
time.

\begin{utadtimeline}{31}
\ganttbar[name=model]{Model Development (v1--v8)}{1}{15} \\
\ganttbar{Validation Sweeps}{13}{15} \\
\ganttmilestone{Final Test Eval}{14} \\
\ganttbar{Data Augmentation Study (v1--v5)}{12}{20} \\
\ganttbar{Self-Directed Upskilling}{13}{31} \\
\ganttbar[name=infer]{Inference Pipeline Setup}{20}{20} \\
\ganttbar[name=deploy, bar/.append style={fill=utadgrey!35, draw=utadgrey!60}]
  {Deployment Demo (planned)}{21}{31} \\
% Finish-to-start dependency arrows between named bars
\ganttlink{model}{infer}
\ganttlink{infer}{deploy}
\end{utadtimeline}

\section{Model Development}

\subsection{Objective}
Develop a convolutional classifier that generalizes well to unseen
Fashion-MNIST images, improving on a naive baseline through a sequence of
individually-measured changes, each validated on the fixed validation split
before being kept.

\subsection{Model Evolution (v1 -- v8)}
\begin{utadchain}
\utadchainbox{start}{v1}{Baseline MLP}
\utadchainbox{mid}{v2}{Simple CNN}
\utadchainbox{mid}{v3}{Normalization}
\utadchainbox{mid}{v4}{Deeper net}
\utadchainbox{mid}{v5}{Dropout}
\utadchainbox{mid}{v6}{Augmentation}
\utadchainbox{mid}{v7}{LR schedule}
\utadchainbox{final}{v8}{Final model}
\end{utadchain}

\begin{itemize}
\item \textbf{v1 -- Baseline MLP.} A single fully-connected hidden layer on
  flattened pixels, as a measured floor for everything that follows.
\item \textbf{v2 -- Simple CNN.} Two convolution-and-pool blocks, which cut
  the validation error sharply by exploiting spatial structure the MLP
  ignored.
\item \textbf{v5 -- Dropout.} Added dropout to the dense head to close the
  train/validation gap that appeared once the network was deep enough to
  overfit.
\item \textbf{v8 -- Final model.} Locked architecture and hyperparameters
  after the validation sweeps; this is the only version measured on the
  test set.
\end{itemize}

\subsection{Technical Details}
\begin{lstlisting}[language=Python]
# One training epoch: forward pass, cross-entropy loss, backprop, SGD step.
def train_epoch(model, loader, optimizer, device):
    model.train()
    running = 0.0
    for images, labels in loader:
        images, labels = images.to(device), labels.to(device)
        optimizer.zero_grad()
        logits = model(images)
        loss = F.cross_entropy(logits, labels)   # softmax + NLL
        loss.backward()
        optimizer.step()
        running += loss.item() * images.size(0)
    return running / len(loader.dataset)
\end{lstlisting}

\subsection{Objective Function}
The network outputs class scores $z_c$ that a softmax turns into a
probability distribution over the $C = 10$ classes:
\begin{equation}
  p_c = \frac{e^{z_c}}{\sum_{j=1}^{C} e^{z_j}}.
\end{equation}
For a single labeled example with one-hot target $y$, training minimizes the
cross-entropy between the target and the predicted distribution:
\begin{equation}
  \mathcal{L} = -\sum_{c=1}^{C} y_c \log p_c .
\end{equation}
Over a training set of $N$ examples, the optimizer minimizes the mean loss,
taking a stochastic-gradient step with learning rate $\eta$ each iteration:
\[
  J(\theta) = \frac{1}{N}\sum_{i=1}^{N} \mathcal{L}_i,
  \qquad
  \theta \leftarrow \theta - \eta\, \nabla_\theta J(\theta).
\]

\begin{utadimportant}[Measured once, not tuned against]
Every architecture and hyperparameter choice above was selected using the
validation split only. The held-out test set was evaluated exactly once, on
the final model (v8), so the reported test accuracy is an honest estimate of
generalization rather than a number the model was tuned toward.
\end{utadimportant}

\subsection{Results}
{\utadtable
\begin{tabularx}{\textwidth}{@{}Xrr@{}}
\tblheadrow \tblhead{Model} & \tblhead{Val.\ accuracy} & \tblhead{Params} \\
v1 -- Baseline MLP & 87.9\% & 0.10M \\
v2 -- Simple CNN & 90.6\% & 0.23M \\
v5 -- + Dropout & 91.8\% & 0.23M \\
v8 -- Final model & 93.1\% & 0.41M \\
\end{tabularx}}

\medskip
\noindent The final model reached \textbf{92.7\%} accuracy on the held-out
test set --- close to its validation figure, indicating little overfitting
to the validation split.

\chapter*{Development Log}
\addcontentsline{toc}{chapter}{Appendix A: Development Log}

\utadweeklog{1 (Week 1)}
  {Set up the data pipeline and reproducible training loop; established the
   validation split; trained the baseline MLP (v1) as a measured floor.}
  {Python, PyTorch, NumPy.}
  {Reproducible pipeline and a measured baseline (v1, 87.9\% val).}

\utadweeklog{2 (Week 2)}
  {Introduced convolutional blocks and batch normalization; iterated the
   model through several revisions (v2--v4); began plotting learning curves.}
  {PyTorch, torchvision, Matplotlib.}
  {Convolutional model clearly beat the MLP baseline; first learning-curve
   plots for review.}

\utadweeklog{3 (Week 3)}
  {Added dropout and data augmentation (v5--v6); ran validation sweeps over
   learning rate and augmentation strength; introduced a learning-rate
   schedule (v7).}
  {PyTorch, Weights \& Biases (logging).}
  {Closed the train/validation gap; selected hyperparameters entirely on the
   validation split.}

\utadweeklog{4 (Week 4)}
  {Locked the final model (v8); ran the single test-set evaluation;
   assembled an inference pipeline and a 24-test regression suite to a green
   result.}
  {PyTorch, Git, GitLab.}
  {Final test accuracy of 92.7\%; green regression suite; reproducible
   end-to-end run.}

\chapter*{References}
\addcontentsline{toc}{chapter}{Appendix B: References}
\begin{utadrefs}
\utadref Goodfellow, I., Bengio, Y., \& Courville, A. \emph{Deep Learning}.
  MIT Press. \url{https://www.deeplearningbook.org/}
\utadref Xiao, H., Rasul, K., \& Vollgraf, R. \emph{Fashion-MNIST: a Novel
  Image Dataset for Benchmarking Machine Learning Algorithms}.
  \url{https://arxiv.org/abs/1708.07747}
\utadref Kingma, D. P., \& Ba, J. \emph{Adam: A Method for Stochastic
  Optimization}. \url{https://arxiv.org/abs/1412.6980}
\utadref PyTorch Team. \emph{PyTorch Documentation}.
  \url{https://pytorch.org/docs/}
\end{utadrefs}

\end{document}
